Because the two generators swap members only within their own blocks, the product action \(S_2\times S_2\) has one two-member orbit in each block and arm. Hence the orbit partition is \[\{O_{1,0},O_{1,1},O_{2,0},O_{2,1}\}.\] If \(q_{\mathrm{tar}}\) is concentrated on \(v_{1,0}\), its arm-\(1\) lift lies in \(O_{1,1}\), so the \((1,O_{1,1})\) coordinate of \(b_{\mathrm{LP}}\) equals one. The mate-control state \((v_{1,1},0)\) lies in \(O_{1,0}\subseteq\mathcal U_0\), has zero mass in the \((1,O_{1,1})\) coordinate, and is not a legal arm-\(1\) selector. Thus the arm-changing analogue cannot meet the stacked target equations.

Separately, if a legal selector has \(q^{\mathrm{obs}}_{w,1}(O_{1,1})=0\), then \(O_{1,1}\) is the sole target-supported arm-\(1\) orbit and the premise of \(\texttt{prop:disjoint-orbit-witness}\) holds. For \(k\le M\), that proposition gives \(\mathcal C_w=\{0\}\), \(\tau_T=K\), and conditional coverage zero. For \(k=M+1\), \(\texttt{prop:rank-boundary}\) gives \(\mathcal C_w=\mathbb R\) and conditional coverage one.